\documentclass{uai2026} % for initial submission
% \documentclass[accepted]{uai2026} % after acceptance

\usepackage[american]{babel}
\usepackage{natbib}
\bibliographystyle{plainnat}
\renewcommand{\bibsection}{\subsubsection*{References}}
\usepackage{mathtools}
\usepackage{booktabs}
\usepackage{graphicx}

\title{ThinkDet: An Experimental VLM Adapter for Functional Query Grounding}

% Author block is automatically hidden for anonymous submission when the
% 'accepted' option is not used.
\author[1]{Anonymous Authors}
\affil[1]{Anonymous Institution}

\begin{document}
\maketitle

\begin{abstract}
ThinkDet is an experimental adapter that injects query-conditioned hidden states
from a frozen vision-language model into a frozen open-vocabulary detector.
The mechanism is lightweight: InternVL hidden states are pooled into eight
summary tokens by a Text Memory Augmenter and fused into selected
GroundingDINO decoder layers. Across the saved repo artifacts, the evidence is
mixed rather than uniformly positive. Historical repo artifacts show small
gains on a held-out affordance benchmark together with improved top-1
calibration, but a corrected rerun with query-conditioned scoring changes that
picture: Stage 2 improves affordance top-1 hit rate only slightly
(\(16.99\%\rightarrow17.77\%\)) and does not preserve the earlier calibration
gain. The main early checkpoints also regress on standard benchmarks,
including official COCO AP and all reported RefCOCO-family splits for the
primary Stage-2 checkpoint. We therefore present ThinkDet as a narrow
functional-query stress-test result, not as a general improvement to
open-vocabulary detection. The main value of the current project is an honest
research artifact: a plausible adapter mechanism, a small affordance-specific
gain, and clear evidence that this gain does not transfer cleanly to standard
detection settings.
\end{abstract}

\section{Introduction}\label{sec:intro}
Open-vocabulary object detectors such as GroundingDINO \citep{groundingdino2023}
work well on standard referential queries, but they can become unstable on
functional prompts such as \emph{something to talk on} or
\emph{something to carry things in}. ThinkDet explores whether a frozen
vision-language model can supply extra semantics for that narrow regime without
retraining the detector end-to-end.

The core mechanism is straightforward. We extract query-conditioned InternVL
hidden states, pool them into a small set of learned summary tokens, and fuse
those tokens into GroundingDINO's decoder text memory. This yields a compact
adapter with a low trainable-parameter budget and a clear interface between the
VLM and detector.

The important repo-level conclusion is not that ThinkDet broadly solves
reasoning-guided detection. The saved artifacts instead show a mixed result:
small gains on a custom functional-query benchmark, alongside regressions on
standard COCO and RefCOCO-family evaluation. This paper draft is therefore
best read as a narrow experimental study rather than a general detector claim.

\paragraph{Contributions.}
\begin{enumerate}
\item We implement a parameter-efficient adapter that injects query-conditioned VLM features into a frozen detector.
\item We define and evaluate a small held-out affordance benchmark with calibration-focused metrics.
\item We show a mixed empirical result: a small affordance-benchmark gain under corrected query-conditioned scoring, but regressions on standard COCO and RefCOCO-family evaluation for the main early checkpoints.
\item We argue that ThinkDet should be framed as a narrow functional-query stress-test study, not a general improvement to open-vocabulary detection.
\end{enumerate}

\section{Related Work}\label{sec:related}
\subsection{Uncertainty in Object Detection}
Uncertainty estimation for deep detectors has been studied using Monte Carlo dropout \citep{gal2016dropout}, deep ensembles \citep{lakshminarayanan2017deep}, and evidential approaches \citep{sensoy2018evidential}. These methods primarily estimate uncertainty after prediction. In contrast, ThinkDet aims to reduce ambiguity-driven uncertainty before prediction by injecting additional semantic information into the detector's text pathway.

\subsection{Open-Vocabulary Object Detection}
Open-vocabulary and grounded detection systems such as GLIP \citep{glip2022}, DetCLIP \citep{detclip2022}, and GroundingDINO \citep{groundingdino2023} extend detection beyond fixed category vocabularies. GroundingDINO combines DINO-style detection \citep{dino2023} with text conditioning and is the detector backbone in our work.

\subsection{Parameter-Efficient Adaptation}
Parameter-efficient fine-tuning methods such as adapters \citep{houlsby2019adapter}, LoRA \citep{lora2022}, and prefix tuning \citep{prefix2021} motivate our design. TMA is closest in spirit to memory/prefix augmentation, but its tokens are input-dependent because they are produced from a live VLM forward pass conditioned on both image and query.

\subsection{VLM-Augmented Detection and Calibration}
Jointly trained multimodal detectors such as GLIP \citep{glip2022} and FIBER \citep{fiber2022} can improve grounding but typically require large-scale retraining. ThinkDet keeps both the detector and VLM frozen and inserts a lightweight bridge. Post-hoc calibration methods such as temperature scaling \citep{guo2017calibration} adjust scores after prediction, whereas ThinkDet modifies internal representations and lets us test whether query-conditioned semantic priors change ranking quality or confidence reliability on ambiguous queries.

\section{Problem Formulation: Uncertainty on Functional Queries}\label{sec:problem}
\subsection{The Uncertainty Gap}
Let a detector \(D\) with text encoder \(E_{\text{text}}\) and visual encoder \(E_{\text{vis}}\) take an image \(I\) and query \(q\), and produce candidate boxes with confidence scores \(\{(b_i, c_i)\}\).

For referential queries, \(E_{\text{text}}(q)\) is often sufficiently discriminative, and confidence tends to concentrate on correct detections. For functional queries, \(E_{\text{text}}(q)\) can be diffuse because the text encoder does not reliably encode affordance-level semantics. This often yields:
\begin{itemize}
\item high predictive dispersion across candidates,
\item poor confidence calibration,
\item and missed detections due to weak ranking.
\end{itemize}

\subsection{Uncertainty Metrics}
We use both currently available and planned uncertainty metrics:
\begin{enumerate}
\item \textbf{Top-1 confidence calibration (measured)}: expected calibration error (ECE), Brier score, and negative log-likelihood computed from stored per-query top-1 scores and correctness labels.
\item \textbf{Top-\(K\) predictive entropy (planned next artifact)}: entropy over normalized top-\(K\) detection scores per query, once the evaluator saves candidate score vectors.
\item \textbf{Accuracy@0.5}: grounding accuracy at intersection-over-union threshold 0.5.
\end{enumerate}

This lets us evaluate not only whether ThinkDet improves accuracy, but also whether accuracy and confidence reliability move together on functional queries.

\section{Method}\label{sec:method}
\subsection{Architecture Overview}
ThinkDet combines a frozen GroundingDINO detector with a frozen InternVL model \citep{internvl2024} and a trainable Text Memory Augmenter.

At inference time:
\begin{enumerate}
\item The image is processed by GroundingDINO and InternVL.
\item The query text is processed by GroundingDINO's BERT encoder and jointly with the image in InternVL.
\item We extract text-conditioned visual hidden states \(H_{\text{vlm}}\) from an intermediate InternVL layer.
\item TMA pools \(H_{\text{vlm}}\) into \(M=8\) summary tokens and concatenates them to the detector text memory at decoder layers \(\{1,3,5\}\).
\end{enumerate}

The injected tokens act as a semantic prior that can improve query disambiguation without modifying the detector backbone or the VLM.

\subsection{Text Memory Augmenter}
Given text-conditioned InternVL visual features \(H_{\text{vlm}} \in \mathbb{R}^{N_{\text{vis}} \times d_{\text{vlm}}}\), TMA projects them into the detector text-memory dimension and pools them with learned queries:
\begin{align}
K &= \mathrm{LN}(W_{\text{proj}} H_{\text{vlm}}), \\
V &= K, \\
Q &= Q_{\text{learned}}, \\
Z &= \mathrm{MHA}(Q, K, V), \\
\text{aug}_l &= \alpha_l Z.
\end{align}

Here \(Q_{\text{learned}} \in \mathbb{R}^{M \times d_{\text{dino}}}\) are learned query vectors (broadcast across the batch), \(Z\) are the pooled summary tokens, and \(\alpha_l\) is a learned layerwise scaling factor at decoder layer \(l\).

\subsection{Layerwise Scaling}
In the current implementation, \(\alpha_l\) is a static scalar per injection layer (not a query-conditioned uncertainty estimator). We therefore interpret \(\alpha_l\) as a layerwise augmentation-strength parameter and treat uncertainty estimates as output-space measurements (calibration metrics now, predictive entropy in a follow-up evaluation pass).

\subsection{Training Protocol}
\paragraph{Stage 1 (COCO semantic pretraining).}
TMA is trained on standard COCO detection (train2017) to learn a generic VLM-to-detector feature translation. In the newer TMA pipeline, only TMA parameters are trainable in Stage 1.

\paragraph{Stage 2 (heads + TMA).}
TMA and selected detection heads are trainable while the detector backbone and VLM remain frozen. The repository contains multiple Stage 2 variants, including COCO-based Stage 2 and a separate RefCOCO-based Stage 2 run.

\paragraph{No affordance supervision.}
The main ThinkDet affordance results are obtained without affordance-specific training labels; improvements on affordance grounding are emergent from COCO-based training and VLM prior injection.

\section{Experiments}\label{sec:experiments}
\subsection{Setup}
\paragraph{Models.}
The baseline used by the current repo artifacts is GroundingDINO Swin-T OGC.
ThinkDet adds InternVL-based TMA with \(M=8\) summary tokens and injection
layers \(\{1,3,5\}\).

\paragraph{Reported checkpoints.}
The affordance results use:
\begin{itemize}
\item a Stage-1 TMA checkpoint trained on COCO (\texttt{layer9\_tma\_m8\_full\_2ep\_7gpu\_20260218\_152649}),
\item and a Stage-2 checkpoint (\texttt{layer9\_tma\_m8\_8gpu\_20260219\_010134}).
\end{itemize}

\paragraph{Benchmarks.}
We report:
\begin{enumerate}
\item a held-out affordance grounding benchmark (512 queries, 8 affordance categories),
\item COCO val2017 sanity-check detection evaluation,
\item and RefCOCO/RefCOCO+/RefCOCOg referring expression evaluation.
\end{enumerate}

\subsection{Main Result: Held-Out Affordance Grounding}\label{sec:aff-main}
\paragraph{Thesis note.}
The affordance numbers below are the corrected rerun produced by the active
query-conditioned scoring path. Older historical repo artifacts reported
larger gains under an earlier max-logit ranking shortcut; those numbers are
useful only as provenance, not as the canonical thesis result.

Table~\ref{tab:affordance-main} reports top-1 hit rate at intersection-over-union
threshold 0.5. Under the corrected scorer, ThinkDet improves overall
performance by only 0.78 percentage points on this benchmark, with the largest
gain on \texttt{carry\_in}. The gains are narrow and non-uniform across
categories, and \texttt{talk\_on} regresses, which is why the result should be
treated as a stress-test outcome rather than as evidence of broad detector
improvement.

\begin{table*}[t]
\centering
\caption{Held-Out Affordance Benchmark (512 Queries): Top-1 Hit Rate at IoU \(\geq 0.5\).}
\label{tab:affordance-main}
\begin{tabular}{lrrrr}
\toprule
Affordance & Baseline & Stage 1 & Stage 2 & \(\Delta\) (Stage2--Base) \\
\midrule
\texttt{carry\_in}   & 15.62\% & 17.19\% & \textbf{20.31\%} & \textbf{+4.69pp} \\
\texttt{cut\_with}   &  6.25\% & \textbf{7.81\%} &  6.25\%  & +0.00pp \\
\texttt{drink\_from} & 37.50\% & \textbf{39.06\%} & \textbf{39.06\%} & +1.56pp \\
\texttt{eat\_with}   &  4.69\% &  4.69\% &  4.69\%         & +0.00pp \\
\texttt{read}        & 18.75\% & 18.75\% & 18.75\%         & +0.00pp \\
\texttt{ride}        & 31.25\% & 31.25\% & \textbf{32.81\%} & +1.56pp \\
\texttt{sit\_on}     &  4.69\% &  4.69\% &  4.69\%         & +0.00pp \\
\texttt{talk\_on}    & \textbf{17.19\%} & \textbf{17.19\%} & 15.62\% & -1.56pp \\
\midrule
Overall (\texttt{hit@0.5\_top1}) & 16.99\% & 17.58\% & \textbf{17.77\%} & \textbf{+0.78pp} \\
\bottomrule
\end{tabular}
\end{table*}

\subsection{Uncertainty Analysis (Current Artifact + Next Measurement)}\label{sec:uncertainty}
The current affordance evaluator stores one top-1 score per query (\texttt{top1\_score}) and the resulting top-1 IoU, which enables calibration-focused analysis. It does not yet store top-\(K\) candidate score vectors, so predictive-entropy analysis is specified as a next step.

\subsubsection{Top-1 Confidence Calibration (Measured)}
Table~\ref{tab:calibration} shows that the corrected scorer rerun does not
preserve the earlier calibration story from historical artifacts. Stage 2
improves top-1 hit rate slightly and AUROC modestly, but ECE, Brier score, and
negative log-likelihood all worsen relative to the baseline. This suggests
that the earlier large calibration gains were at least partly tied to the
older ranking shortcut rather than the adapter alone.

\begin{table}[t]
\centering
\caption{Calibration-Focused Uncertainty Metrics on the Held-Out Affordance Benchmark (512 Queries).}
\label{tab:calibration}
\begin{tabular}{lrrr}
\toprule
Metric & Baseline & Stage 1 & Stage 2 \\
\midrule
\texttt{hit@0.5\_top1} & 0.1699 & 0.1758 & \textbf{0.1777} \\
ECE (10 bins) & \textbf{0.0632} & 0.0744 & 0.0773 \\
Brier score & \textbf{0.1429} & 0.1486 & 0.1498 \\
Negative log-likelihood & \textbf{0.4654} & 0.4836 & 0.4854 \\
AUROC (score \(\rightarrow\) correctness) & 0.5882 & 0.5730 & \textbf{0.5997} \\
\bottomrule
\end{tabular}
\end{table}

Mean top-1 confidence also drops slightly from 0.1067 (baseline) to 0.1004
(Stage 2) while accuracy increases from 0.1699 to 0.1777, but the bin-based
calibration metrics worsen. Lower raw confidence alone should therefore not be
interpreted as improved calibration on this benchmark.

\subsubsection{Predictive Entropy (Next Artifact)}
To measure predictive dispersion directly, the affordance evaluator should save top-\(K\) scores or logits for each query (for example, \(K=50\)). This enables predictive entropy, entropy reduction by affordance category, and risk-coverage analysis using richer uncertainty signals than top-1 confidence.

\subsection{COCO Sanity Check}
Table~\ref{tab:coco} reports a clean COCO val2017 evaluation for a Stage-1 checkpoint. ThinkDet regresses on standard category detection in this setting, which is consistent with our claim that the method targets functional-query uncertainty rather than general-purpose COCO detection gains.

\begin{table}[t]
\centering
\caption{COCO val2017 Clean Evaluation Sanity Check (4952 Images).}
\label{tab:coco}
\begin{tabular}{lrrrr}
\toprule
Model & mAP & mAP\(_{50}\) & mAP\(_{75}\) & Trainable Params \\
\midrule
Baseline (clean eval) & \textbf{43.23} & \textbf{56.99} & \textbf{47.12} & --- \\
ThinkDet (Stage 1, TMA only) & 39.66 & 50.84 & 43.40 & 1.58M (0.13\%) \\
\bottomrule
\end{tabular}
\end{table}

\subsection{RefCOCO / RefCOCO+ / RefCOCOg}
Table~\ref{tab:refcoco} reports top-1 REC accuracy at IoU \(\geq 0.5\) (confidence threshold 0.05) for a Stage-2 checkpoint evaluated across eight standard splits. The baseline wins all splits, supporting the claim that ThinkDet is not a general REC improvement module and is most useful in the functional-query regime.

\begin{table*}[t]
\centering
\caption{RefCOCO Family Evaluation: Top-1 Accuracy at IoU \(\geq 0.5\) (Confidence Threshold 0.05).}
\label{tab:refcoco}
\begin{tabular}{lrrr}
\toprule
Dataset / Split & Baseline & ThinkDet (Stage 2) & \(\Delta\) \\
\midrule
refcoco val   & \textbf{50.62} & 48.61 & -2.01pp \\
refcoco testA & \textbf{56.43} & 51.83 & -4.60pp \\
refcoco testB & \textbf{45.53} & 45.22 & -0.31pp \\
refcoco+ val   & \textbf{51.33} & 47.18 & -4.15pp \\
refcoco+ testA & \textbf{55.82} & 48.83 & -6.99pp \\
refcoco+ testB & \textbf{46.23} & 46.00 & -0.23pp \\
refcocog val  & \textbf{60.50} & 57.58 & -2.92pp \\
refcocog test & \textbf{60.50} & 57.72 & -2.78pp \\
\bottomrule
\end{tabular}
\end{table*}

\section{Analysis and Ablations}\label{sec:analysis}
\subsection{Layerwise Scaling as Augmentation Strength}
For the Stage-2 checkpoint used in the affordance benchmark, the learned layerwise scaling parameters are approximately:
\begin{center}
\begin{tabular}{lr}
\toprule
Injection Layer & Learned \(\alpha\) \\
\midrule
Layer 1 & \(\approx 0.0885\) \\
Layer 3 & \(\approx 0.0966\) \\
Layer 5 & \(\approx 0.0983\) \\
\bottomrule
\end{tabular}
\end{center}

These values are static per layer and should not be interpreted as per-query uncertainty estimates. A better future analysis is to correlate query difficulty with query-dependent effective augmentation statistics, such as \(\lVert \alpha_l Z_l \rVert\) or attention mass onto injected tokens.

\subsection{Planned Ablations for the Submission Version}
The following ablations are part of our current paper plan and remain to be finalized with artifact-backed uncertainty metrics:
\begin{itemize}
\item mixed-query-mode comparison for Stage 1 pretraining,
\item injection layer ablation \(\{1\}, \{3\}, \{5\}, \{1,3\}, \{1,3,5\}\),
\item number of summary tokens \(M\),
\item and InternVL source-layer selection.
\end{itemize}

\section{Discussion}\label{sec:discussion}
Under the corrected scorer, ThinkDet improves affordance grounding only
slightly on average while degrading performance on standard REC and COCO
settings. The calibration picture is mixed rather than uniformly improved.
This pattern supports a narrow stress-test interpretation rather than a
universal accuracy-improvement claim. In particular:
\begin{itemize}
\item functional queries appear more uncertainty-limited on the text side, so VLM-derived priors can help;
\item referential queries are often already handled well by the baseline, so extra tokens may introduce noise;
\item standard category detection is not primarily limited by the same semantic ambiguity, so augmentation can perturb an already strong solution.
\end{itemize}

We can interpret TMA as injecting a prior-like signal into the detector's posterior over boxes. In this view, InternVL features provide a prior over likely objects conditioned on image and query context, while the detector combines this with its existing visual-text matching machinery.

\subsection{Limitations}
\begin{enumerate}
\item We do not optimize an explicit uncertainty objective.
\item The layerwise scaling parameter is static per layer and not query-conditioned.
\item The current affordance benchmark is modest in size (512 queries).
\item The current artifact stores only top-1 scores, not top-\(K\) candidate score vectors needed for predictive-entropy analysis.
\item We only evaluate one VLM backbone family in the current draft.
\end{enumerate}

\section{Conclusion}\label{sec:conclusion}
We presented ThinkDet, a parameter-efficient method for injecting semantic
priors from a frozen vision-language model into a frozen detector for
ambiguity-heavy grounding. On the corrected held-out affordance rerun,
ThinkDet improves grounding accuracy only slightly, does not improve the main
top-1 calibration metrics, and regresses on standard COCO and RefCOCO-family
benchmarks. These results support a narrow interpretation of ThinkDet as an
experimental functional-query adapter and analysis artifact, not as a general
open-vocabulary detection improvement.

\bibliography{thinkdet_uai2026}

\newpage
\onecolumn
\appendix

\section{Training Details}\label{app:train}
Table~\ref{tab:train-details} summarizes the training hyperparameters used in the paper draft. Some values are run-dependent because the repository contains older and newer pipelines.

\begin{table}[h]
\centering
\caption{Training Hyperparameters (Paper Draft Summary).}
\label{tab:train-details}
\begin{tabular}{lcc}
\toprule
Hyperparameter & Stage 1 & Stage 2 \\
\midrule
GPUs & up to 8 A100 & up to 8 A100 \\
Batch size per GPU & 2 & 2 \\
Gradient accumulation & 4 & 4 \\
Effective batch size & 64 & 64 \\
Learning rate (TMA) & \(3\times 10^{-5}\) to run-specific & \(3\times 10^{-5}\) to run-specific \\
Learning rate (heads) & stage-dependent & stage-dependent \\
Warmup ratio & 0.05 & 0.05 \\
Reported affordance checkpoints & epoch 2 & epoch 2 \\
AMP & fp16 & fp16 \\
\(\alpha\) init & 0.3 & loaded from Stage 1 \\
\(\alpha\) floor & optional, script-dependent & optional, script-dependent \\
\bottomrule
\end{tabular}
\end{table}

\section{Held-Out Affordance Benchmark Construction}\label{app:benchmark}
The held-out benchmark contains 512 queries split evenly across eight affordance categories. Each sample contains a natural-language functional description paired with a COCO-style image containing at least one valid target.

\begin{table}[h]
\centering
\caption{Held-Out Affordance Categories (64 Queries Each).}
\label{tab:aff-types}
\begin{tabular}{ll}
\toprule
Affordance Type & Example Query \\
\midrule
\texttt{talk\_on} & ``something to talk on'' \\
\texttt{carry\_in} & ``something to carry things in'' \\
\texttt{sit\_on} & ``something to sit on'' \\
\texttt{cut\_with} & ``something to cut with'' \\
\texttt{drink\_from} & ``something to drink from'' \\
\texttt{eat\_with} & ``something to eat with'' \\
\texttt{read} & ``something to read'' \\
\texttt{ride} & ``something to ride'' \\
\bottomrule
\end{tabular}
\end{table}

\section{Uncertainty Metrics Details}\label{app:uncertainty}
\subsection{Current Artifact Metrics}
The current evaluator stores \texttt{top1\_score} and \texttt{top1\_iou} for each query. This supports:
\begin{itemize}
\item ECE with fixed confidence bins,
\item Brier score,
\item negative log-likelihood,
\item and correctness-prediction AUROC.
\end{itemize}

\subsection{Planned Predictive Entropy}
Once top-\(K\) candidate scores are logged, we compute entropy from normalized scores:
\begin{equation}
H(p) = - \sum_{i=1}^{K} \hat{p}_i \log \hat{p}_i,
\end{equation}
where \(\hat{p}_i = \operatorname{softmax}(c_i)\) over top-\(K\) candidate logits \(c_i\).

\section{Artifact Notes on Stage Naming}\label{app:stages}
The repository contains multiple Stage 2 variants:
\begin{itemize}
\item a COCO-based Stage 2 (\texttt{train\_stage2\_tma.py}), and
\item a separate RefCOCO joint-training Stage 2 (\texttt{train\_refcoco\_stage2\_tma.py}) that trains on RefCOCO, RefCOCO+, and RefCOCOg.
\end{itemize}
Both variants train TMA and detector heads while keeping the detector backbone and VLM frozen.

\end{document}
